\documentclass{article}
\usepackage[T1]{fontenc}
\usepackage{lmodern}
\usepackage{graphicx}
\usepackage{amsmath,amssymb}
\usepackage[a4paper,margin=2cm]{geometry}
\usepackage[absolute,overlay]{textpos}
\setlength{\TPHorizModule}{1mm}
\setlength{\TPVertModule}{1mm}
\usepackage[indonesian]{babel}
\usepackage{hyperref}
\setlength{\emergencystretch}{2em}
% unit: Halaman 23

\begin{document}

% === Halaman 23 (llm) ===

\section*{RC6}

\begin{itemize}
    \item RC6 juga dibuat oleh Ronald Rivest, Matt Robshaw, Ray Sidney, dan Yiqun Lisa Yin dari Laboratorium RSA (Rivest, 1998).
    \item RC6 merupakan pengembangan RC5. RC6 dibuat untuk mengikuti kompetisi Advanvced Encryption Standard atau AES (akan dijelaskan pada materi selanjutnya).
    \vspace{10mm}
    \item RC6 melakukan enkripsi terhadap blok data berukuran 128 bit, sedangkan ukuran kuncinya bervariasi 128, 192, 256 sampai 2040 bit. Cara pembangkitan kunci internalnya identik dengan RC5.
    \vspace{10mm}
    \item Secara struktur, RC6 mirip dengan RC5, begitu juga operasi-operasinya seperti XOR, penjumlahan, dan rotasi. RC6 mempunyai operasi tambahan yaitu perkalian.
    \item RC6 dapat dipandang sebagai dua buah RC5 yang dioperasikan secara parallel.
\end{itemize}

\end{document}
